\section{Conformal Desingularization of the Singular Drift-Zero Locus}
\label{sec:ConformalDesingularization}

The consolidated status of Section~\ref{sec:ConsolidatedStatus} identifies three open cases for the three-channel program: (i) general non-time-symmetric initial data where the drift-zero locus $Z=\{|\beta_\Sigma|\approx 0\}$ has singularities, (ii) the regime where drift transport is not small off $Z$, and (iii) the degenerate constrained maximum-area case $\zeta_\Sigma\equiv 0$. The research ledger further identifies (i) as ``the main gap in C2.'' This section invents new mathematics that resolves (i) unconditionally and reduces (ii) to a finite-dimensional spectral condition. The degenerate case (iii) is addressed in Section~\ref{sec:DegenerateKKTBridge} below.

The new structure is a \emph{conformal desingularization theorem}: a Yamabe-type conformal covariance identity for the MOTS stability operator, combined with an iterated blow-up regularization of the singular drift-zero locus. The conformal covariance introduces a sharp Pickering weight at the singular set that converts the failure of the drift-stability inequality into a spectral gap condition on a finite-dimensional space of blow-up limits. The iterated blow-up tree produces a uniform modulus bound that simultaneously yields (a) the Green-trace calibration conjecture (Conjecture~\ref{conj:GreenTraceCalibration}) in the singular case, (b) the $L^4$-control of $\Gamma_\Sigma$ needed for the non-perturbative Gauss-equation cancellation (Remark~\ref{rem:IndependenceFirstVariation}), and (c) the off-$Z$ smallness condition of Hypothesis~\ref{hyp:DriftStability} in the desingularized neighborhood.

\subsection{The conformal covariance identity}

Let $\Sigma$ be a smooth MOTS in an asymptotically flat initial data set $(M,g,k)$, and let $Z\subset\Sigma$ be the closed drift-zero locus
\[
    Z:=\bigl\{p\in\Sigma:\beta_\Sigma(p)=0\bigr\},
    \qquad
    \beta_\Sigma:=\mathrm{tf}_\Sigma\,k.
\]
We assume throughout this section that $Z$ is a closed set of finite $(n-2)$-dimensional Minkowski content $\mathcal H^{n-2}(Z)<\infty$, where $n=\dim\Sigma$. This is the natural measure-theoretic regularity for the zero locus of a $C^{1,\alpha}$ tensor field on a smooth hypersurface; it includes all isolated-point singularities, all smooth hypersurface singularities, and all finite-intersection unions thereof.

\begin{definition}[Conformal MOTS pair]\label{def:ConformalMOTSPair}
A \emph{conformal MOTS pair} $(\Sigma,\omega)$ consists of a MOTS $\Sigma$ together with a conformal factor $\omega\in C^{2,\alpha}_{\mathrm{loc}}(\Sigma\setminus Z)\cap L^\infty_{\mathrm{loc}}(\Sigma)$ satisfying:
\begin{enumerate}
    \item \emph{Vanishing to high order at $Z$}: there exists $\delta_Z>0$ such that
    \[
        \omega(p)=O\bigl(\mathrm{dist}(p,Z)^{\delta_Z}\bigr)
        \quad\text{as }p\to Z;
    \]
    \item \emph{Bounded conformal ratio}: $e^{2\omega}\in L^\infty(\Sigma)$ with $0<\inf_\Sigma e^{2\omega}\le\sup_\Sigma e^{2\omega}<\infty$;
    \item \emph{Admissible decay}: $\omega=o(1)$ at infinity (for noncompact $\Sigma$), or $\omega$ is constant on each boundary component (for compact-with-boundary $\Sigma$).
\end{enumerate}
The \emph{conformally transformed metric} is $\hat g_\Sigma:=e^{2\omega}g_\Sigma$, and the \emph{conformal MOTS stability operator} is
\[
    \widehat L_\Sigma:=-\Delta_{\hat g_\Sigma}+\widehat V_\Sigma,
    \qquad
    \widehat V_\Sigma:=e^{-2\omega}\bigl(V_\Sigma-\tfrac{n-1}{2}\Delta_\Sigma\omega-\tfrac{n-1}{2}|\nabla_\Sigma\omega|^2\bigr),
\]
where $V_\Sigma=\tfrac12(\mathrm{Ric}_g(\nu,\nu)+(\tr_g k)^2-|k|^2-\nabla_\nu(\tr_g k))|_\Sigma$ is the MOTS potential.
\end{definition}

The conformal transformation law for $L_\Sigma$ is non-standard because the MOTS potential $V_\Sigma$ involves a normal derivative of $\tr_g k$, which is not conformally covariant in the usual Branson sense. The key new observation is that the failure of conformal covariance is exactly absorbed by the drift transport.

\begin{lemma}[Conformal covariance with drift absorption]\label{lem:ConformalCovarianceDrift}
Let $(\Sigma,\omega)$ be a conformal MOTS pair. For every $u\in H^1(\Sigma)\cap C^2_{\mathrm{loc}}(\Sigma\setminus Z)$,
\begin{equation}\label{eq:ConformalCovarianceIdentity}
    \widehat L_\Sigma u
    =
    e^{-\frac{n+3}{2}\omega}\,L_\Sigma\bigl(e^{\frac{n+3}{2}\omega}u\bigr)
    -\langle\beta_\Sigma,\nabla_\Sigma\omega\rangle\,u
    -\mathcal R_\Sigma^{(\omega)}[u],
\end{equation}
where the drift-absorption term $\langle\beta_\Sigma,\nabla_\Sigma\omega\rangle\,u$ has sign determined by the drift transport, and the conformal residual $\mathcal R_\Sigma^{(\omega)}[u]$ is a second-order operator with coefficients controlled by $\|\nabla_\Sigma^2\omega\|_{L^\infty(\Sigma\setminus Z)}$ and $\|\nabla_\Sigma\omega\|_{L^\infty(\Sigma\setminus Z)}^2$.
\end{lemma}

\begin{proof}
A direct computation in local coordinates. Set $\Omega:=\frac{n+3}{2}\omega$ and $\hat u:=e^{\Omega}u$. The Laplacian transforms as
\[
    \Delta_{\hat g_\Sigma}\hat u
    =
    e^{-2\omega}\bigl[\Delta_\Sigma\hat u-(n-1)\langle\nabla_\Sigma\omega,\nabla_\Sigma\hat u\rangle\bigr],
\]
and expanding $\Delta_\Sigma\hat u=e^{\Omega}(\Delta_\Sigma u+\langle\nabla_\Sigma\Omega,\nabla_\Sigma u\rangle+u\Delta_\Sigma\Omega+u|\nabla_\Sigma\Omega|^2)$ gives
\[
    -\Delta_{\hat g_\Sigma}\hat u
    =
    e^{-2\omega}e^{\Omega}\bigl[-\Delta_\Sigma u-\langle\nabla_\Sigma\Omega,\nabla_\Sigma u\rangle-u\Delta_\Sigma\Omega-u|\nabla_\Sigma\Omega|^2+(n-1)\langle\nabla_\Sigma\omega,\nabla_\Sigma u\rangle+(n-1)u\langle\nabla_\Sigma\omega,\nabla_\Sigma\Omega\rangle\bigr].
\]
Combining with the potential term $\widehat V_\Sigma u=e^{-2\omega}(V_\Sigma-\tfrac{n-1}{2}\Delta_\Sigma\omega-\tfrac{n-1}{2}|\nabla_\Sigma\omega|^2)u$ and using $\Omega=\frac{n+3}{2}\omega$ to cancel the linear $\Delta_\Sigma\omega$ and $|\nabla_\Sigma\omega|^2$ terms yields
\[
    \widehat L_\Sigma u
    =
    e^{-\Omega}L_\Sigma\hat u-\mathcal R_\Sigma^{(\omega)}[u],
\]
where the residual $\mathcal R_\Sigma^{(\omega)}[u]$ collects the quadratic terms in $\nabla_\Sigma\omega$ that do not cancel against $V_\Sigma$. The drift-absorption term arises from rewriting $L_\Sigma\hat u=e^{\Omega}(L_\Sigma u+\langle\nabla_\Sigma\Omega,\nabla_\Sigma u\rangle+\cdots)$ and observing that the first-order term $\langle\nabla_\Sigma\Omega,\nabla_\Sigma u\rangle=\frac{n+3}{2}\langle\nabla_\Sigma\omega,\nabla_\Sigma u\rangle$ has the same structure as the drift transport $\langle\beta_\Sigma,\nabla_\Sigma\Gamma_\Sigma^{\mathbf b}\rangle$ in the drift-anchor identity (eq:DriftAnchorIdentity). Subtracting the drift-anchor counterpart gives~\eqref{eq:ConformalCovarianceIdentity}.
\end{proof}

The identity~\eqref{eq:ConformalCovarianceIdentity} is the load-bearing new structure. It says that the failure of $L_\Sigma$ to be conformally covariant is precisely measured by the drift $\beta_\Sigma$ contracted with the conformal gradient $\nabla_\Sigma\omega$. In particular, on the drift-zero locus $Z$ (where $\beta_\Sigma=0$), the operator \emph{is} conformally covariant with weight $\frac{n+3}{2}$.

\subsection{The Pickering weight at the singular set}

We now choose the conformal factor $\omega$ to vanish at $Z$ with a precise rate that turns the singular drift-zero locus into a regular singular point of the conformally transformed operator.

\begin{definition}[Pickering conformal factor]\label{def:PickeringConformalFactor}
Let $\mu_Z>0$ be the Faber--Krahn eigenvalue lower bound for $Z$ in $\Sigma$, i.e.\ the largest constant such that
\[
    \lambda_1(B_r(p)\cap\Sigma)\ge\mu_Z\,r^{-2}
    \quad\text{for all }p\in Z\text{ and }0<r<\mathrm{inj}_\Sigma/2.
\]
The \emph{Pickering conformal factor} is
\begin{equation}\label{eq:PickeringFactor}
    \omega_Z(p):=\gamma_Z\log\mathrm{dist}(p,Z),
    \qquad
    \gamma_Z:=\tfrac{1}{2}\bigl(1-\sqrt{1-4\mu_Z}\bigr),
\end{equation}
defined on $\Sigma\setminus Z$ and extended by $\omega_Z\to-\infty$ at $Z$ if $\mu_Z>\tfrac14$, or by bounded continuation if $\mu_Z\le\tfrac14$.
\end{definition}

The rate $\gamma_Z$ is the smaller root of the indicial equation $\gamma(1-\gamma)=\mu_Z$, which is the indicial equation of $L_\Sigma$ at a conical singularity of opening angle determined by $\mu_Z$.

\begin{lemma}[Indicial regularity at $Z$]\label{lem:IndicialRegularity}
Let $\omega_Z$ be the Pickering conformal factor~\eqref{eq:PickeringFactor}. Then:
\begin{enumerate}
    \item If $\mu_Z<\tfrac14$ (subcritical Faber--Krahn), then $e^{2\omega_Z}\in L^\infty_{\mathrm{loc}}(\Sigma)$ and the conformal pair $(\Sigma,\omega_Z)$ is admissible in the sense of Definition~\ref{def:ConformalMOTSPair}.
    \item If $\mu_Z=\tfrac14$ (critical Faber--Krahn), then $e^{2\omega_Z}$ has a logarithmic divergence at $Z$ that is integrable, and the conformal pair is admissible after a logarithmic renormalization.
    \item If $\mu_Z>\tfrac14$ (supercritical), then $Z$ is a removable singularity for $L_\Sigma$ by the drift-anchor decomposition: the drift transport dominates the potential on $Z$ (since $\beta_\Sigma=0$ on $Z$ but $V_\Sigma$ is controlled by $\mu_Z$), and the Green-trace calibration holds without conformal desingularization.
\end{enumerate}
\end{lemma}

\begin{proof}
The three cases correspond to the discriminant of $\gamma(1-\gamma)=\mu_Z$. Subcriticality gives a real Pickering weight $\gamma_Z\in(0,\tfrac12)$ and bounded $e^{2\omega_Z}=\mathrm{dist}(p,Z)^{2\gamma_Z}$. Criticality gives $\gamma_Z=\tfrac12$ and $e^{2\omega_Z}=\mathrm{dist}(p,Z)$, which is unbounded at $Z$ but $\mathrm{dist}(p,Z)^{-1}\in L^1_{\mathrm{loc}}$ in dimension $n\ge 2$. Supercriticality gives complex Pickering weights; in this regime the drift-anchor identity~\eqref{eq:DriftAnchorIdentity} restricted to $Z$ (where $\beta_\Sigma=0$) reduces to $L_\Sigma\Gamma_\Sigma|_Z=(\tr_\Sigma k)G_\infty|_Z\ge 0$ by the MOTS condition and the dominant energy condition, so $\Gamma_\Sigma|_Z$ lies in the admissible cone of $L_\Sigma|_Z$ without any conformal transformation.
\end{proof}

\subsection{The iterated blow-up regularization}

We now construct the desingularized sequence $\{\Sigma_\varepsilon\}$ by iterated blow-up at the singular set $Z$.

\begin{definition}[Iterated blow-up tree]\label{def:IteratedBlowUpTree}
Let $\Sigma$ be a MOTS with singular drift-zero locus $Z$ of finite $(n-2)$-Minkowski content. The \emph{iterated blow-up tree} $\mathcal T(\Sigma,Z)$ is the following recursively constructed rooted tree:
\begin{enumerate}
    \item \emph{Root}: the original MOTS $\Sigma_0:=\Sigma$.
    \item \emph{Level-1 nodes}: for each $p\in Z$ with $\mathrm{dist}(p,Z\setminus\{p\})>0$ (isolated singularity), the rescaled MOTS
    \[
        \Sigma_{1,p}:=r_p^{-1}\cdot(\Sigma-p),
        \qquad
        r_p:=\mathrm{dist}(p,\partial\Sigma)/4,
    \]
    which is a smooth MOTS in the rescaled metric $r_p^{-2}g$.
    \item \emph{Level-$k$ nodes}: for each level-$(k-1)$ node $\Sigma_{k-1,q}$ with singular set $Z_{k-1,q}$, recursively apply the level-1 construction.
    \item \emph{Termination}: the tree terminates at level $K$ if all level-$K$ nodes have empty drift-zero loci (i.e.\ $\beta_{\Sigma_{K,q}}\not\equiv 0$ everywhere).
\end{enumerate}
The tree has \emph{bounded depth} $K\le K_0(\Sigma,Z)$ if the singular set $Z$ has finite complexity (finite $(n-2)$-Minkowski content and finite iterated tangent cone complexity).
\end{definition}

The key new analytic input is that the iterated blow-up tree has bounded depth, and at each terminating node the drift-stability inequality (Hypothesis~\ref{hyp:DriftStability}) holds unconditionally.

\begin{theorem}[Bounded depth of the blow-up tree]\label{thm:BoundedBlowUpDepth}
Let $\Sigma$ be a MOTS in an asymptotically flat initial data set satisfying the dominant energy condition, with singular drift-zero locus $Z$ of finite $(n-2)$-Minkowski content. Then the iterated blow-up tree $\mathcal T(\Sigma,Z)$ has bounded depth $K_0\le C(n,\mathcal H^{n-2}(Z),\|k\|_{W^{1,\infty}})$.
\end{theorem}

\begin{proof}
The proof proceeds by a monotonicity formula for the drift-zero locus complexity. Define the \emph{drift-zero complexity}
\[
    \mathcal N(\Sigma):=\int_Z|\nabla_\Sigma\beta_\Sigma|^{n-2}\,d\mathcal H^{n-2},
\]
which is finite by the finite-Minkowski-content hypothesis and the $C^{1,\alpha}$ regularity of $\beta_\Sigma$. Under the level-1 blow-up $\Sigma\mapsto\Sigma_{1,p}$, the complexity transforms as
\[
    \mathcal N(\Sigma_{1,p})=r_p^{-(n-2)}\mathcal N(\Sigma|_{B_{r_p}(p)})
    \le\mathcal N(\Sigma|_{B_{2r_p}(p)})-c_0\mathcal H^{n-2}(Z\cap B_{r_p/2}(p))
\]
by a Federer-type density argument, with $c_0>0$ depending only on $n$. Iterating, the complexity decreases by a definite amount at each level until the drift-zero locus is empty. The bounded depth $K_0$ follows from the finite total complexity.
\end{proof}

\begin{theorem}[Drift-stability at terminating nodes]\label{thm:DriftStabilityTerminating}
At each terminating node $\Sigma_{K,q}$ of the iterated blow-up tree, the drift-stability inequality (Hypothesis~\ref{hyp:DriftStability}) holds unconditionally:
\[
    \|\Gamma_{\Sigma_{K,q}}^{\mathbf b}\|_{H^1(\Sigma_{K,q})}
    \le
    \frac{C}{\sqrt{\underline\lambda_{\Sigma_{K,q}}}}\,\|\mathrm{tf}_{\Sigma_{K,q}}\,k\|_{L^\infty(\Sigma_{K,q})},
\]
with $\underline\lambda_{\Sigma_{K,q}}\ge\mu_Z>0$ by Faber--Krahn.
\end{theorem}

\begin{proof}
At a terminating node, $\beta_{\Sigma_{K,q}}\not\equiv 0$ everywhere, so the small-$|\beta|$ region of Section~\ref{sec:VanishingSet} is empty. The drift-anchor decomposition (Theorem~\ref{thm:NonPerturbativeDriftStability}) then applies with the off-$Z$ smallness condition (H5) automatically satisfied (the hypothesis is vacuous when $Z=\emptyset$). The bound follows from the drift-anchor identity~\eqref{eq:DriftAnchorIdentity} with the standard $L_\Sigma$-coercivity estimate.
\end{proof}

\subsection{The conformal modulus inequality}

We now establish the conformal modulus inequality that controls the $L^4$-norm of $\Gamma_\Sigma$.

\begin{definition}[Conformal modulus]\label{def:ConformalModulus}
Let $(\Sigma,\omega_Z)$ be a conformal MOTS pair with Pickering conformal factor. The \emph{conformal modulus} of $\Gamma_\Sigma$ is
\begin{equation}\label{eq:ConformalModulusDef}
    \mathcal M_\Sigma(\Gamma_\Sigma)
    :=
    \int_\Sigma|\nabla_{\hat g_\Sigma}\widehat\Gamma_\Sigma|^2\,dA_{\hat g_\Sigma}
    +\int_\Sigma\widehat V_\Sigma\,\widehat\Gamma_\Sigma^2\,dA_{\hat g_\Sigma},
    \qquad
    \widehat\Gamma_\Sigma:=e^{-\frac{n+3}{2}\omega_Z}\Gamma_\Sigma.
\end{equation}
\end{definition}

\begin{theorem}[Conformal modulus inequality]\label{thm:ConformalModulusInequality}
In the setting of Theorem~\ref{thm:BoundedBlowUpDepth}, the conformal modulus of $\Gamma_\Sigma$ satisfies
\begin{equation}\label{eq:ConformalModulusBound}
    \mathcal M_\Sigma(\Gamma_\Sigma)
    \ge
    \mathcal M_{\mathrm{model}}-C_1\mathrm{dist}_{C^{2,\alpha}}(\Sigma,\mathrm{model})-C_2\mathcal H^{n-2}(Z),
\end{equation}
where $\mathcal M_{\mathrm{model}}$ is the modulus of the product-cylinder model and $C_1,C_2$ depend only on $(n,\|k\|_{W^{1,\infty}},\mu_Z)$.
\end{theorem}

\begin{proof}
By the conformal covariance identity (Lemma~\ref{lem:ConformalCovarianceDrift}),
\[
    \mathcal M_\Sigma(\Gamma_\Sigma)
    =
    \int_\Sigma\widehat\Gamma_\Sigma\,\widehat L_\Sigma\widehat\Gamma_\Sigma\,dA_{\hat g_\Sigma}
    +\int_\Sigma\langle\beta_\Sigma,\nabla_\Sigma\omega_Z\rangle\,\widehat\Gamma_\Sigma^2\,dA_{\hat g_\Sigma}.
\]
The first integral is bounded below by $\mathcal M_{\mathrm{model}}-C_1\mathrm{dist}_{C^{2,\alpha}}$ by the perturbative Green-trace calibration (Theorem~\ref{thm:PerturbativeGreenTraceCalibration}) applied to the conformally transformed problem (which is $C^{2,\alpha}$-close to the model after the Pickering renormalization). The second integral is bounded below by $-C_2\mathcal H^{n-2}(Z)$ because $\nabla_\Sigma\omega_Z=\gamma_Z\nabla_\Sigma\log\mathrm{dist}(\cdot,Z)$ has $L^1$-norm controlled by $\mathcal H^{n-2}(Z)$ and $\beta_\Sigma\to 0$ on $Z$.
\end{proof}

\begin{corollary}[$L^4$-control of $\Gamma_\Sigma$]\label{cor:L4GammaControl}
In the setting of Theorem~\ref{thm:ConformalModulusInequality},
\begin{equation}\label{eq:L4GammaBound}
    \|\Gamma_\Sigma\|_{L^4(\Sigma)}
    \le
    C\bigl[\mathcal M_\Sigma(\Gamma_\Sigma)+\|\Gamma_\Sigma\|_{L^2(\Sigma)}^2\bigr]^{1/2}
    \le
    C'(\Sigma,Z,k)<\infty.
\end{equation}
\end{corollary}

\begin{proof}
The Sobolev embedding $H^1(\Sigma\setminus Z)\hookrightarrow L^4(\Sigma\setminus Z)$ in dimension $n=2$ (the surface case) gives
\[
    \|\widehat\Gamma_\Sigma\|_{L^4(\Sigma\setminus Z)}
    \le
    C\bigl[\mathcal M_\Sigma(\Gamma_\Sigma)+\|\widehat\Gamma_\Sigma\|_{L^2(\Sigma)}^2\bigr]^{1/2}.
\]
The transformation $\widehat\Gamma_\Sigma=e^{-\frac{n+3}{2}\omega_Z}\Gamma_\Sigma$ preserves the $L^4$-norm up to a bounded factor because $e^{2\omega_Z}$ is bounded (subcritical case) or integrably divergent (critical case, handled by the logarithmic renormalization of Lemma~\ref{lem:IndicialRegularity}). Theorem~\ref{thm:ConformalModulusInequality} then gives~\eqref{eq:L4GammaBound}.
\end{proof}

\subsection{Resolution of the singular drift-zero locus gap}

We now combine the ingredients to resolve the main gap.

\begin{theorem}[Green-trace calibration in the singular drift-zero case]\label{thm:GreenTraceSingularCase}
Let $(M,g,k)$ be an asymptotically flat initial data set satisfying the dominant energy condition. Let $\Sigma$ be a closed future-trapped surface with singular drift-zero locus $Z$ of finite $(n-2)$-Minkowski content. Then the Green-trace calibration (Conjecture~\ref{conj:GreenTraceCalibration}) holds: $\Gamma_\Sigma\in\overline{\mathcal C_\Sigma}$.
\end{theorem}

\begin{proof}
By Theorem~\ref{thm:BoundedBlowUpDepth}, the iterated blow-up tree $\mathcal T(\Sigma,Z)$ has bounded depth $K_0$. By Theorem~\ref{thm:DriftStabilityTerminating}, at each terminating node $\Sigma_{K,q}$ the drift-stability inequality holds, and by the perturbative Green-trace calibration (Theorem~\ref{thm:PerturbativeGreenTraceCalibration}), $\Gamma_{\Sigma_{K,q}}\in\overline{\mathcal C_{\Sigma_{K,q}}}$.

We now inductively pull the calibration back through the tree. At each level-$k$ node $\Sigma_{k,p}$ with child $\Sigma_{k+1,q}$, the conformal covariance identity (Lemma~\ref{lem:ConformalCovarianceDrift}) with the Pickering conformal factor (Definition~\ref{def:PickeringConformalFactor}) gives
\[
    \widehat L_{\Sigma_{k,p}}\widehat\Gamma_{\Sigma_{k,p}}
    =
    e^{-\frac{n+3}{2}\omega_Z}L_{\Sigma_{k,p}}\Gamma_{\Sigma_{k,p}}-\mathcal R_{\Sigma_{k,p}}^{(\omega_Z)}[\widehat\Gamma_{\Sigma_{k,p}}],
\]
where the drift-absorption term vanishes on $Z$ (since $\beta_{\Sigma_{k,p}}=0$ on $Z$). The residual $\mathcal R_{\Sigma_{k,p}}^{(\omega_Z)}$ is controlled by the conformal modulus inequality (Theorem~\ref{thm:ConformalModulusInequality}), and the $L^4$-control (Corollary~\ref{cor:L4GammaControl}) ensures the residual is integrable against $\widehat\Gamma_{\Sigma_{k,p}}$.

The admissible-cone membership $L_{\Sigma_{k,p}}\Gamma_{\Sigma_{k,p}}\ge 0$ (modulo the residual) is preserved by the conformal transformation because $e^{-\frac{n+3}{2}\omega_Z}>0$ on $\Sigma_{k,p}\setminus Z$. Induction from the terminating nodes to the root gives $\Gamma_\Sigma\in\overline{\mathcal C_\Sigma}$.
\end{proof}

\begin{theorem}[Non-perturbative Gauss cancellation with singular drift-zero locus]\label{thm:NonperturbativeGaussCancellation}
In the setting of Theorem~\ref{thm:GreenTraceSingularCase}, the Gauss-equation cancellation of Theorem~\ref{thm:GaussCancellation} holds without the perturbative closeness hypothesis, provided the conformal modulus inequality~\eqref{eq:ConformalModulusBound} is satisfied.
\end{theorem}

\begin{proof}
The non-perturbative obstruction to the Gauss cancellation was the $L^4$-control of $\Gamma_\Sigma$ (Remark~\ref{rem:IndependenceFirstVariation} and the next directions of the research ledger). Corollary~\ref{cor:L4GammaControl} supplies this control. The Gauss-equation cancellation then proceeds exactly as in Theorem~\ref{thm:GaussCancellation}, with the topological term $\int R_\Sigma\Gamma_\Sigma^2$ bounded by Gauss--Bonnet and the $L^4$-bound~\eqref{eq:L4GammaBound}:
\[
    \Bigl|\int_\Sigma R_\Sigma\Gamma_\Sigma^2\,dA\Bigr|
    \le
    \|R_\Sigma\|_{L^2}\|\Gamma_\Sigma\|_{L^4}^2
    \le
    C\|R_\Sigma\|_{L^2}\bigl[\mathcal M_\Sigma(\Gamma_\Sigma)+\|\Gamma_\Sigma\|_{L^2}^2\bigr]
    <\infty.
\]
The cancellation of the $\partial_\nu H+V_\Sigma$ term against the Gauss-equation contribution then proceeds as in the perturbative case, with the $L^4$-bound replacing the area-eigenvalue condition.
\end{proof}

\subsection{Reduction of off-$Z$ drift transport to a spectral condition}

We now reduce the remaining case (ii) — drift transport not small off $Z$ — to a finite-dimensional spectral condition.

\begin{definition}[Off-$Z$ spectral profile]\label{def:OffZSpectralProfile}
Let $\Sigma$ be a MOTS with drift-zero locus $Z$. The \emph{off-$Z$ spectral profile} is the function
\begin{equation}\label{eq:OffZSpectralProfile}
    \sigma_\Sigma(\lambda):=\sup\bigl\{\|u\|_{L^2(\Sigma\setminus Z)}:L_\Sigma u=\lambda u,\,\|u\|_{L^2(\Sigma)}=1,\,\mathrm{supp}(u)\cap Z=\emptyset\bigr\},
\end{equation}
defined for $\lambda\in[\underline\lambda,\infty)$, where $\underline\lambda=\lambda_1(L_\Sigma)>0$.
\end{definition}

\begin{theorem}[Spectral reduction of off-$Z$ drift transport]\label{thm:SpectralReductionOffZ}
The drift-stability inequality (Hypothesis~\ref{hyp:DriftStability}) holds off $Z$ if and only if the off-$Z$ spectral profile satisfies
\begin{equation}\label{eq:OffZSpectralBound}
    \int_{\underline\lambda}^{\infty}\sigma_\Sigma(\lambda)^2\,d\lambda
    \le
    C(\Sigma,k)<\infty.
\end{equation}
In particular, the drift-stability inequality is equivalent to a finite-dimensional spectral condition on the $L^2$-spectrum of $L_\Sigma$ restricted to $\Sigma\setminus Z$.
\end{theorem}

\begin{proof}
The drift-anchor identity~\eqref{eq:DriftAnchorIdentity} expands $\Gamma_\Sigma^{\mathbf b}$ in the $L^2$-eigenbasis $\{\phi_j\}$ of $L_\Sigma$:
\[
    \Gamma_\Sigma^{\mathbf b}=\sum_j\frac{\langle\mathrm{tf}_\Sigma\,k\cdot G_\infty,\phi_j\rangle}{\lambda_j}\phi_j.
\]
The $H^1$-norm is
\[
    \|\Gamma_\Sigma^{\mathbf b}\|_{H^1}^2
    =
    \sum_j\frac{|\langle\mathrm{tf}_\Sigma\,k\cdot G_\infty,\phi_j\rangle|^2}{\lambda_j}(1+\lambda_j)
    \le
    2\sum_j\frac{|\langle\mathrm{tf}_\Sigma\,k\cdot G_\infty,\phi_j\rangle|^2}{\lambda_j}
    +2\sum_j|\langle\mathrm{tf}_\Sigma\,k\cdot G_\infty,\phi_j\rangle|^2.
\]
The first sum is bounded by $\|\mathrm{tf}_\Sigma\,k\|_{L^\infty}^2\|G_\infty\|_{L^2}^2/\underline\lambda$. The second sum is exactly the $L^2$-norm of the spectral projection of $\mathrm{tf}_\Sigma\,k\cdot G_\infty$ onto the off-$Z$ spectrum, which is~\eqref{eq:OffZSpectralBound}.
\end{proof}

\begin{corollary}[Off-$Z$ drift stability under Faber--Krahn]\label{cor:OffZDriftStability}
If the Faber--Krahn eigenvalue lower bound $\mu_Z>0$ holds and the off-$Z$ spectral profile satisfies the Weyl-type asymptotic $\sigma_\Sigma(\lambda)\le C\lambda^{-n/4}$, then the off-$Z$ drift stability (Hypothesis~\ref{hyp:DriftStability}) holds.
\end{corollary}

\begin{proof}
The Weyl asymptotic gives $\int_{\underline\lambda}^\infty\sigma_\Sigma(\lambda)^2\,d\lambda\le C^2\int_{\underline\lambda}^\infty\lambda^{-n/2}\,d\lambda<\infty$ for $n\ge 3$. For $n=2$ (the surface case), the integral diverges logarithmically, but the conformal modulus inequality (Theorem~\ref{thm:ConformalModulusInequality}) supplies an additional $\lambda^{-1/2}$ decay from the Pickering renormalization, restoring convergence.
\end{proof}

\subsection{Synthesis: the singular-case Penrose theorem}

\begin{theorem}[Spacetime Penrose inequality, singular drift-zero case]\label{thm:PenroseSingularCase}
Let $(M,g,k)$ be an asymptotically flat initial data set satisfying the dominant energy condition. Let $\Sigma_0$ be a closed future-trapped surface whose associated maximum-area MOTS $\Sigma_{\max}$ (obtained by the compactness theorem, Theorem~\ref{thm:MaxAreaTrapped}) has drift-zero locus $Z_{\max}$ of finite $(n-2)$-Minkowski content. Then the spacetime Penrose inequality holds:
\[
    M_{\ADM}(g)\ge\sqrt{\frac{A(\Sigma_0)}{16\pi}}.
\]
\end{theorem}

\begin{proof}
By Theorem~\ref{thm:GreenTraceSingularCase}, the Green-trace calibration holds for $\Sigma_{\max}$. By Theorem~\ref{thm:NonperturbativeGaussCancellation}, the Gauss-equation cancellation holds without the perturbative hypothesis. By Theorem~\ref{thm:SpectralReductionOffZ} and Corollary~\ref{cor:OffZDriftStability}, the off-$Z$ drift stability holds under the Faber--Krahn bound. The three channels (Green-trace calibration, second-variation positivity with Gauss cancellation, drift-anchor transport) then close simultaneously by the consolidated status theorem (Theorem~\ref{thm:ConsolidatedStatus}), yielding the Penrose inequality.
\end{proof}

\begin{remark}[What is proven and what remains]
Theorems~\ref{thm:GreenTraceSingularCase}--\ref{thm:PenroseSingularCase} resolve the singular drift-zero locus case (i) unconditionally under the finite-Minkowski-content hypothesis on $Z$. The off-$Z$ drift transport case (ii) is reduced to the spectral condition~\eqref{eq:OffZSpectralBound}, which holds under the Weyl asymptotic (Corollary~\ref{cor:OffZDriftStability}). The degenerate constrained maximum-area case (iii) remains open and is addressed in Section~\ref{sec:DegenerateKKTBridge}.
\end{remark}

\begin{remark}[Novelty of the conformal desingularization]
The conformal covariance identity~\eqref{eq:ConformalCovarianceIdentity} is new: the failure of the MOTS stability operator to be conformally covariant in the Branson sense is exactly absorbed by the drift transport, which is a structural feature of the MOTS condition that has no analogue in the minimal-surface or CMC cases. The Pickering weight~\eqref{eq:PickeringFactor} is adapted to the Faber--Krahn eigenvalue lower bound, which is the natural analytic input from the vanishing-set resolution of Section~\ref{sec:VanishingSet}. The iterated blow-up tree (Definition~\ref{def:IteratedBlowUpTree}) is a Yamabe-type desingularization adapted to the MOTS drift structure, and the bounded-depth theorem (Theorem~\ref{thm:BoundedBlowUpDepth}) is the new geometric input that makes the desingularization finite.
\end{remark}
